% State-specific — Financial Literacy: Budgeting, Saving, and Investing
\section{Financial Literacy --- Budgeting, Saving, and Investing}

\begin{learningGoals}
\begin{itemize}
    \item Create and analyze a simple budget
    \item Compare saving and investing options
\end{itemize}
\end{learningGoals}

\begin{conceptBox}{Budgeting Basics}
A \textbf{budget} is a plan for your money.

\textbf{Income} $-$ \textbf{Expenses} $=$ \textbf{Savings}

\begin{itemize}[leftmargin=*]
    \item \textbf{Fixed expenses} stay the same each month (rent, phone plan).
    \item \textbf{Variable expenses} change (food, entertainment).
    \item \textbf{Investing} means putting money into something (stocks, bonds) hoping it grows over time. It has \textbf{risk} --- you could lose money too.
\end{itemize}
\end{conceptBox}

\begin{workedExample}{Making a Monthly Budget}
Monthly income: $\$1{,}200$. Fixed expenses: $\$700$. Variable expenses: $\$350$.

Savings $= 1{,}200 - 700 - 350 = \$150$

Savings rate: $\frac{150}{1{,}200} = 0.125 = 12.5\%$
\answerHighlight{$\$150$ saved, which is $12.5\%$ of income.}
\end{workedExample}

\mascotSays{Experts suggest saving at least $10\%$ of your income. Even small amounts add up!}

\begin{practiceBox}{Budgeting and Saving}
\resetProblems

\prob Income: $\$2{,}000$/month. Fixed: $\$1{,}100$. Variable: $\$600$. How much is left to save?
\answerExplain{$\$300$}{$2{,}000 - 1{,}100 - 600 = \$300$.}

\prob You save $\$50$/month at $3\%$ simple interest for $1$ year. How much interest do you earn on the total saved?
\answerExplain{$\$9$}{Total saved: $50 \times 12 = \$600$. Average balance over a year is about $\$300$, but simple interest on $\$600$ for half a year averages out. Approximate: $I = 600 \times 0.03 \times 0.5 = \$9$.}

\prob You want to save $\$1{,}200$ in $8$ months. How much per month?
\answerExplain{$\$150$/month}{Divide: $1{,}200 \div 8 = \$150$ per month.}

\prob Your food budget is $\$400$/month. You spend $\$425$ one month. By what percent did you go over?
\answerExplain{$6.25\%$}{Over by $\$25$. Percent: $25 \div 400 = 0.0625 = 6.25\%$.}

\prob A savings account earns $2\%$/year. An investment averages $7\%$/year but could lose money. Name one reason to choose each.
\answerExplain{Savings: safe; Investment: higher growth}{Savings accounts are low-risk --- your money is safe. Investments can grow faster but carry the risk of losing value.}

\end{practiceBox}
